% Aturan format; tidak perlu diubah saat hanya mengisi laporan.
\usepackage[left=4cm,right=3cm,top=3cm,bottom=3cm]{geometry}
\usepackage{fontspec}
% Gunakan LuaLaTeX atau XeLaTeX, bukan pdfLaTeX.
\IfFontExistsTF{Times New Roman}
  {\setmainfont{Times New Roman}}
  {\setmainfont{TeX Gyre Termes}}
\IfFontExistsTF{Consolas}
  {\setmonofont[Scale=0.9]{Consolas}}
  {}
\usepackage{amsmath,amssymb,graphicx,booktabs,array,tabularx,longtable}
% Tambahan untuk buku manual: diagram, grafik, kode, tautan panjang.
\usepackage{tikz,pgfplots,fancyvrb,xcolor}
\usetikzlibrary{arrows.meta,positioning,calc}
\pgfplotsset{compat=1.18}
\usepackage{xurl}
\usepackage{setspace,indentfirst,titlesec,caption,float,placeins,enumitem}
% BibTeX menangani references.bib; gaya sitasi angka sesuai urutan kutipan.
\usepackage[backend=biber,style=numeric,sorting=none]{biblatex}
\addbibresource{references.bib}
\usepackage[hidelinks]{hyperref}
\onehalfspacing
\setlength{\parindent}{1cm}
\setlength{\parskip}{0pt}
\setlength{\emergencystretch}{2em}
\raggedbottom
\pagestyle{plain}
\setlist{leftmargin=0.8cm,topsep=6pt,itemsep=3pt,parsep=0pt}
\renewcommand{\chaptername}{BAB}
\renewcommand{\contentsname}{DAFTAR ISI}
\renewcommand{\listfigurename}{DAFTAR GAMBAR}
\renewcommand{\listtablename}{DAFTAR TABEL}
\renewcommand{\figurename}{Gambar}
\renewcommand{\tablename}{Tabel}
\renewcommand{\bibname}{DAFTAR PUSTAKA}
\renewcommand{\thechapter}{\Roman{chapter}}
\renewcommand{\thesection}{\arabic{chapter}.\arabic{section}}
\renewcommand{\thesubsection}{\thesection.\arabic{subsection}}
\renewcommand{\thefigure}{\arabic{chapter}.\arabic{figure}}
\renewcommand{\thetable}{\arabic{chapter}.\arabic{table}}
\renewcommand{\theequation}{\arabic{chapter}.\arabic{equation}}
\setcounter{tocdepth}{2}
\titleformat{\chapter}[display]
  {\normalfont\bfseries\centering\singlespacing\fontsize{14}{18}\selectfont}
  {BAB \thechapter}{5pt}{}
\titleformat{\section}{\normalfont\normalsize\bfseries}{\thesection}{0.7em}{}
\titleformat{\subsection}{\normalfont\normalsize\bfseries}{\thesubsection}{0.7em}{}
\titlespacing{\chapter}{0pt}{0pt}{24pt}
\titlespacing{\section}{0pt}{14pt}{6pt}
\titlespacing{\subsection}{0pt}{10pt}{4pt}
\DeclareCaptionFont{ukcaption}{\fontsize{10}{12}\selectfont}
\captionsetup{font={ukcaption,singlespacing},skip=6pt}
\graphicspath{{figures/}}
% Kode dan perintah: \kode{...} untuk teks pendek (garis bawah ditulis \_),
% lingkungan Perintah untuk blok perintah yang disalin apa adanya.
\newcommand{\kode}[1]{\texttt{#1}}
\DefineVerbatimEnvironment{Perintah}{Verbatim}
  {fontsize=\small,frame=single,framesep=2mm,rulecolor=\color{black!35},baselinestretch=1}
\newcolumntype{Y}{>{\raggedright\arraybackslash}X}
\newcolumntype{L}[1]{>{\raggedright\arraybackslash}p{#1}}
\setlength{\LTpre}{6pt}
\setlength{\LTpost}{6pt}
\definecolor{abu}{HTML}{8A8F98}
\definecolor{aksen}{HTML}{2A78D6}
\definecolor{garis}{HTML}{B9C0C9}

